\documentclass[11pt]{letter}
\usepackage[utf8]{inputenc}
\usepackage[T1]{fontenc}
\usepackage{geometry}
\geometry{margin=1in}

\signature{Francisco Parra-O.\\
Departamento de Ingenier\'ia Inform\'atica\\
Universidad de Santiago de Chile\\
francisco.parra.o@usach.cl}

\address{Francisco Parra-O.\\
Departamento de Ingenier\'ia Inform\'atica\\
Universidad de Santiago de Chile\\
Av. V\'ictor Jara 3659\\
Santiago 9170124, Chile}

\begin{document}

\begin{letter}{Editor-in-Chief\\
Computers \& Geosciences\\
Elsevier}

\opening{Dear Editor,}

I am pleased to submit my manuscript entitled ``\textbf{nowcast: a forcing-agnostic
Rust engine for dynamic geohazard nowcasting, and why forcing resolution, not the
model, sets the skill ceiling}'' for consideration for publication in
\textit{Computers \& Geosciences}.

Susceptibility mapping is mature but static; empirical rainfall thresholds say
\emph{when} but not \emph{where}. Operational early-warning needs both, yet existing
systems weld a specific forcing, trigger and susceptibility model into one
product-specific chain. I present \textbf{nowcast}, an open, dependency-light Rust
engine that expresses time-varying hazard as a susceptibility field times a
composable trigger behind a single interchangeable forcing interface, and I then
use it as an experimental instrument to ask a question a fixed chain cannot: holding
the model fixed, how much does the \emph{resolution of the forcing} determine
attainable skill?

\textbf{Key contributions:}
\begin{itemize}
    \item A forcing-agnostic engine (ten Rust crates) in which rainfall, routed
    discharge, snowmelt and InSAR deformation enter through one interface, with a
    composable multi-trigger and exact, closed-form attribution of every alert.
    \item A backtesting framework with metrics appropriate for rare, spatially
    sparse, month-dated inventories (event-centred contingency, ROC-AUC,
    probability-of-detection-at-area, split-sample calibration).
    \item A reproducible \emph{resolution diagnosis}: on the R\'io Maipo, distributing
    daily forcing and weighting by a real machine-learning susceptibility raster does
    not lift discrimination above chance (AUC $\approx 0.48$), and applying the full
    ensemble of 35 published intensity--duration thresholds does not rescue it; a
    controlled experiment with known ground truth shows the same engine discriminates
    planted sub-daily bursts almost perfectly but loses all operational catch rate
    under daily aggregation. The binding constraint is the forcing resolution, not
    the model.
    \item A validated counterpart: half-hourly satellite forcing pins the
    threshold crossing hours ahead of documented flows across three climatically
    distinct events, and a discharge path validated against observed streamflow on
    the R\'io Itata discriminates floods out of sample (ROC-AUC 0.90).
\end{itemize}

\textbf{Relevance to Computers \& Geosciences:} the work contributes a novel
computational architecture and a methods result for geohazard early-warning, with
emphasis on reproducibility: the engine, its backtesting framework and every
result are open-source and accompanied by runnable examples and data-extraction
scripts. This fits the journal's focus on computational methods and software for
the geosciences.

The manuscript reports an honest negative result (the daily null) alongside its
positive ones; I believe this transparency is a strength for a methods paper and is
well suited to the journal's readership of practitioners and method developers.

This manuscript has not been published elsewhere and is not under consideration by
any other journal. I declare no competing interests. The code and data supporting
this study are publicly available as indicated in the manuscript.

Thank you for considering my submission. I look forward to your response.

\closing{Sincerely,}

\end{letter}
\end{document}
